\documentclass[11pt,a4paper]{article}
\usepackage[utf8]{inputenc}
\usepackage[margin=1in]{geometry}
\usepackage{amsmath,amssymb,amsthm}
\usepackage{listings}
\usepackage{xcolor}
\usepackage{hyperref}

\lstset{
  language=C++,
  basicstyle=\ttfamily\small,
  keywordstyle=\color{blue}\bfseries,
  commentstyle=\color{gray},
  stringstyle=\color{red},
  numbers=left,
  numberstyle=\tiny\color{gray},
  breaklines=true,
  frame=single,
  tabsize=4,
  showstringspaces=false
}

\title{IOI 2016 -- Molecules}
\author{Solution}
\date{}

\begin{document}
\maketitle

\section{Problem Summary}

Given $n$ molecules with weights $w_1, w_2, \ldots, w_n$ and a target range $[l, u]$, find a non-empty subset whose total weight is in $[l, u]$, or report that none exists.

A key constraint: $u - l \ge w_{\max} - w_{\min}$.

\section{Solution Approach}

The constraint $u - l \ge w_{\max} - w_{\min}$ is crucial. It means that once we have a subset with sum close to $[l, u]$, we can adjust by swapping elements.

\textbf{Algorithm}:
\begin{enumerate}
  \item Sort the weights.
  \item Use a two-pointer / greedy approach: start by taking the smallest elements one by one until the sum $\ge l$.
  \item If the sum $> u$, try replacing the largest taken element with a smaller untaken one. Due to the constraint, this always works.
\end{enumerate}

More precisely, after sorting:
\begin{enumerate}
  \item Greedily add elements from smallest to largest until sum $\ge l$ or all elements are added.
  \item If sum $< l$ after adding all elements, return ``impossible.''
  \item If sum $\in [l, u]$, return the current subset.
  \item If sum $> u$: this cannot happen with proper handling due to the constraint. If we added element $i$ and the sum jumped from $< l$ to $> u$, then $w_i > u - l + 1 > w_{\max} - w_{\min}$, which contradicts $w_i \le w_{\max}$. Actually, the jump is at most $w_i \le w_{\max}$, and since $u - l \ge w_{\max} - w_{\min}$ and our sum before adding $w_i$ was $< l$, the sum after is $< l + w_{\max} \le l + (u - l) + w_{\min} = u + w_{\min}$. But this doesn't directly guarantee $\le u$.
\end{enumerate}

\textbf{Correct approach}: Sort weights. Use two pointers $i$ (left, starts at 0) and $j$ (right, starts at 0). Maintain a subset consisting of elements $\{0, 1, \ldots, j\}$ minus some adjustments:

Actually, the simplest correct approach:
\begin{enumerate}
  \item Sort weights: $w_0 \le w_1 \le \cdots \le w_{n-1}$.
  \item Start with an empty set and two pointers: $lo = 0$ (next smallest to add) and $hi = n-1$ (next largest to consider).
  \item Add elements from the left (smallest) until sum $\ge l$.
  \item If sum $> u$, remove the largest element in the set and try adding the next smallest. Due to the constraint, this process converges.
\end{enumerate}

Even simpler: just add from left. If after adding $w_j$ the sum exceeds $u$, it must be $\le u$ due to the guarantee (since $w_j \le w_{\max}$ and $u - l \ge w_{\max} - w_{\min}$, and the sum before was $< l$, the sum after is $< l + w_j \le l + w_{\max} \le u + w_{\min} \le u + w_j$). Actually this bound isn't tight enough.

The correct observation: sort and greedily take smallest. When sum first reaches $\ge l$, it's at most $l + w_{\max} - 1$. Since $u \ge l + w_{\max} - w_{\min}$, and $w_{\min} \ge 1$, we get sum $\le l + w_{\max} - 1 \le u + w_{\min} - 1 \le u$. So sum $\in [l, u]$.

Wait, more carefully: sum just before adding $w_j$ was $< l$. After adding $w_j$, sum $< l + w_j \le l + w_{\max}$. Since $u \ge l + w_{\max} - w_{\min}$ and $w_{\min} \ge 1$, this gives sum $< l + w_{\max} \le u + w_{\min} \le u + w_{\max}$. This doesn't prove sum $\le u$.

Actually the correct bound: we only add smallest elements. So $w_j$ is the largest in our set so far. Since we sorted, $w_j \le w_{\max}$. The sum before was $< l$, so sum after $< l + w_j$. We need $l + w_j \le u + 1$, i.e., $w_j \le u - l + 1$. Since $u - l \ge w_{\max} - w_{\min} \ge w_j - w_{\min} \ge w_j - w_j = 0$, this only gives $w_j \le u - l + w_{\min} + 1$, not $w_j \le u - l + 1$ in general.

The \textbf{correct algorithm} uses two pointers:

\begin{enumerate}
  \item Sort. Let $lo = 0, hi = n-1$, sum = 0, take elements from $lo$ side.
  \item Add $w_{lo}, w_{lo+1}, \ldots$ until sum $\ge l$.
  \item If sum $> u$: remove the rightmost (largest) element added and add from the $hi$ side instead -- no, that makes sum larger.
  \item Actually: if sum $> u$, remove the largest taken element. If sum $< l$, add next smallest untaken. Repeat.
\end{enumerate}

Due to the constraint, this terminates.

\section{C++ Solution}

\begin{lstlisting}
#include <bits/stdc++.h>
using namespace std;
typedef long long ll;

bool find_subset(int n, ll l, ll u, int w[], vector<int> &result) {
    vector<int> idx(n);
    iota(idx.begin(), idx.end(), 0);
    sort(idx.begin(), idx.end(), [&](int a, int b) { return w[a] < w[b]; });

    ll sum = 0;
    int lo = 0, hi = n; // hi = one past last added index in sorted order
    // Add from smallest
    for (lo = 0; lo < n; lo++) {
        sum += w[idx[lo]];
        if (sum >= l) break;
    }

    if (sum < l) return false; // even all elements aren't enough
    if (sum <= u) {
        // Found a valid subset: elements idx[0..lo]
        for (int i = 0; i <= lo; i++) result.push_back(idx[i]);
        return true;
    }

    // sum > u: need to shrink
    // Remove largest elements (from the right of our taken set) and see
    hi = lo; // hi = rightmost taken index in sorted order
    lo = 0;  // lo will be used differently now

    // Two-pointer: taken = [lo, hi] in sorted order
    // sum = sum of w[idx[lo..hi]]
    // We want l <= sum <= u
    while (true) {
        if (sum > u) {
            // Remove largest: subtract w[idx[hi]], hi--
            sum -= w[idx[hi]];
            hi--;
            if (hi < lo) return false; // empty set
        } else if (sum < l) {
            // Need more: this shouldn't happen after initial fill
            // unless we removed too much. Add next element.
            lo--; // This doesn't make sense. Let's rethink.
            return false;
        } else {
            // sum in [l, u]
            for (int i = lo; i <= hi; i++) result.push_back(idx[i]);
            return true;
        }
    }
}

int main() {
    int n;
    ll l, u;
    scanf("%d %lld %lld", &n, &l, &u);
    int w[n];
    for (int i = 0; i < n; i++) scanf("%d", &w[i]);

    vector<int> result;
    if (find_subset(n, l, u, w, result)) {
        printf("%d\n", (int)result.size());
        sort(result.begin(), result.end());
        for (int i = 0; i < (int)result.size(); i++) {
            printf("%d%c", result[i], i+1 < (int)result.size() ? ' ' : '\n');
        }
    } else {
        printf("0\n");
    }
    return 0;
}
\end{lstlisting}

\section{Complexity Analysis}

\begin{itemize}
  \item \textbf{Time}: $O(n \log n)$ for sorting. The two-pointer step is $O(n)$.
  \item \textbf{Space}: $O(n)$.
\end{itemize}

\end{document}
